\chapter{Wstęp}

Tematem niniejszej pracy magisterskiej jest \emph{Wykrywanie Fake News
	przy użyciu metod zespołowego uczenia maszynowego}. Problem \textbf{rozpoznawania fałszywych nowości} jest jednym z
najpoważniejszych wyzwań współczesnego społeczeństwa informacyjnego.

Motywacja do podjęcia tego tematu wynika z realnych zagrożeń, jakie
niesie za sobą dezinformacja oraz potrzeby opracowania skutecznych
narzędzi do jej zwalczania. W ostatnich latach zaobserwowano liczne
przypadki wpływu \textbf{fałszywych nowości} na wyniki wyborów, podziały
społeczne czy reakcje na kryzysy zdrowotne, takie jak pandemia COVID-19.
Informacje nieprawdziwe, ale celowo rozpowszechniane, potrafią wywoływać
panikę, dezinformować w kluczowych sprawach publicznych, a nawet
prowadzić do eskalacji konfliktów międzynarodowych. W tym kontekście
rola zaawansowanych technologii informatycznych w przeciwdziałaniu tym
zjawiskom staje się kluczowa.

Dodatkową motywacją do podjęcia tematu jest rosnące zainteresowanie
zastosowaniem sztucznej inteligencji w analizie danych oraz
ciągła potrzeba udoskonalania istniejących rozwiązań. Współczesne
systemy detekcji \textbf{fałszywych nowości} często bazują na standardowych
algorytmach, które, mimo wysokiej popularności, nie zawsze radzą sobie w
sytuacjach wymagających analizy danych o dużej zmienności, nierównowadze
klas czy subtelnych różnicach semantycznych. Opracowanie nowych metod,
które potrafią efektywnie radzić sobie z tymi wyzwaniami, może znacząco
przyczynić się do zwiększenia skuteczności detekcji \textbf{fałszywych nowości}.

W odpowiedzi na te wyzwania, dziedzina sztucznej inteligencji, a w
szczególności uczenie maszynowe, oferuje zaawansowane narzędzia
umożliwiające skuteczne wykrywanie i klasyfikację dezinformacji. Wśród
tych narzędzi szczególną rolę odgrywają metody zespołowego
\textbf{uczenia maszynowego}, które dzięki
integracji wyników wielu modeli bazowych, pozwalają na osiąganie wyższej
dokładności predykcji i większej stabilności wyników. \textbf{Agregacja metodą bootstrap (bagging)}, \textbf{uczenie wzmacniające (boosting)} czy \textbf{warstwowe łączenie modeli (stacking)} są znane dzięki swojej skuteczności w rozwiązywaniu problemów
nierównowagi klas oraz w zadaniach wymagających analizy dużych i
zróżnicowanych zbiorów danych.

\textbf{Celem pracy} jest kompleksowa analiza tradycyjnych technik uczenia maszynowego oraz nowoczesnych metod zespołowych w kontekście wykrywania \textbf{fałszywych nowości}, ze szczególnym uwzględnieniem badania semantycznych reprezentacji tekstu, wpływu różnych strategii preprocessingu na jakość klasyfikacji oraz identyfikacji najskuteczniejszych podejść w ramach metod zespołowych. Praca ma na celu zbadanie, jak autorskie rozwiązania mogą zwiększyć efektywność systemów detekcji dezinformacji, co pozwala lepiej zrozumieć mechanizmy poprawy skuteczności identyfikacji \textbf{fałszywych nowości}. Na podstawie przeprowadzonych badań planuje się ustalić kluczowe zależności i warunki niezbędne do stworzenia skutecznej procedury wykrywania \textbf{fałszywych nowości}, wzbogacając tym samym wiedzę teoretyczną oraz dostarczając praktyczne narzędzia dla badań naukowych.

Główne założenia pracy obejmują:

\begin{enumerate}
	\def\labelenumi{\arabic{enumi}.}
	\item
	      Opracowanie nowych algorytmów zespołowych wykorzystujących zaawansowane modele semantyczne oraz techniki uczenia do reprezentacji danych tekstowych.
	\item
	      Przeprowadzenie systematycznej analizy efektywności autorskich rozwiązań w kontekście kluczowych miar wydajności klasyfikacyjnej: \textbf{dokładności}, \textbf{mairy F1}, \textbf{obszaru pod krzywą ROC} oraz \textbf{czasu wykonania}.
	\item
	      Praktyczne zastosowanie i walidację proponowanych metod na różnorodnych zbiorach danych zawierających fałszywe i prawdziwe nowości, w tym na uznanych benchmarkowych zestawach, takich jak \textbf{ISOT Fake News Dataset}, \textbf{WELFake} oraz \textbf{BuzzFeed}.
\end{enumerate}

